\documentclass[10pt,a4paper]{article}
\usepackage[margin=1.5cm]{geometry}
\usepackage{enumitem}
\usepackage{titlesec}
\usepackage{parskip}

\titlespacing*{\section}{0pt}{4pt}{2pt}
\titlespacing*{\subsection}{0pt}{3pt}{1pt}
\setlist{nosep, leftmargin=*}
\pagestyle{empty}

\begin{document}

\begin{center}
    \textbf{\large Topic 1: Threat Modelling and What is Security?}
\end{center}

\section*{Formal Definition of Security}
A system is \textbf{secure} if it behaves according to its specification under the presence of an adversary. More precisely: a system is secure if no attacker within a defined threat model can violate the security policy.

\section*{CIA Model}
\begin{itemize}
    \item \textbf{Confidentiality}: Information is only accessible to authorised parties.
    \item \textbf{Integrity}: Information and systems can only be modified by authorised parties in authorised ways.
    \item \textbf{Availability}: Systems and information are accessible to authorised parties when needed.
\end{itemize}

\section*{Lampson's Triad: Policy / Mechanism / Assurance}
\begin{itemize}
    \item \textbf{Policy}: \emph{What} security properties must hold (e.g., ``only admins may delete records'').
    \item \textbf{Mechanism}: \emph{How} the policy is enforced (e.g., access control checks, encryption).
    \item \textbf{Assurance}: \emph{Why} we believe the mechanism correctly enforces the policy (e.g., formal proofs, testing, audits).
\end{itemize}

\section*{STRIDE Threat Categories}
\begin{itemize}
    \item \textbf{S}poofing -- violates \emph{authentication} (pretending to be another entity).
    \item \textbf{T}ampering -- violates \emph{integrity} (unauthorised modification of data).
    \item \textbf{R}epudiation -- violates \emph{non-repudiation} (denying having performed an action).
    \item \textbf{I}nformation Disclosure -- violates \emph{confidentiality} (exposing data to unauthorised parties).
    \item \textbf{D}enial of Service -- violates \emph{availability} (making a service unavailable).
    \item \textbf{E}levation of Privilege -- violates \emph{authorisation} (gaining higher access than allowed).
\end{itemize}

\section*{Threat Modelling Process}
\begin{enumerate}
    \item Identify assets and security objectives.
    \item Create an architecture overview (Data Flow Diagrams with trust boundaries).
    \item Decompose the system into components.
    \item Identify threats using STRIDE per element.
    \item Rate threats (e.g., DREAD) and determine mitigations.
\end{enumerate}

\section*{Example System: Simple Online Store}
\textbf{Components}: Web browser (client), Web server, Database, Payment gateway.\\
\textbf{Trust boundaries}: Internet $\leftrightarrow$ Web server; Web server $\leftrightarrow$ Database; Web server $\leftrightarrow$ Payment gateway.

\subsection*{Threat Analysis (STRIDE applied)}
\begin{tabular}{|l|l|l|}
\hline
\textbf{Threat} & \textbf{Example} & \textbf{Mitigation} \\ \hline
Spoofing & Attacker impersonates a user & Multi-factor authentication \\ \hline
Tampering & Modify order price in transit & TLS + server-side validation \\ \hline
Repudiation & User denies placing an order & Logging + digital signatures \\ \hline
Info.\ Disclosure & SQL injection leaks user data & Input validation, parameterised queries \\ \hline
Denial of Service & Flooding the web server & Rate limiting, CDN \\ \hline
Elevation of Priv. & Normal user accesses admin panel & Role-based access control (RBAC) \\ \hline
\end{tabular}

\section*{Relating to Other Topics}
The formal security definition directly applies to \textbf{web security}: a web application is secure if no attacker (within the web threat model -- e.g., network attacker, script injection) can violate the security policy (e.g., same-origin policy, access control rules). STRIDE maps neatly onto web vulnerabilities: XSS $\to$ Information Disclosure, CSRF $\to$ Spoofing/Tampering, etc.

\end{document}
